% Sources cited in the grammar, including negation, morphology, and phonology.
% Plain LaTeX keeps the existing build independent of a bibliography processor.
\begingroup
\renewcommand{\bibname}{References}
\begin{thebibliography}{Embick \& Noyer 2007}

\bibitem[Baker 1985]{baker1985}
Baker, Mark. 1985.
The Mirror Principle and morphosyntactic explanation.
\textit{Linguistic Inquiry} 16(3): 373--415.
\href{https://brucehayes.org/205/Readings/Baker1985MirrorPrinciple.pdf}{Article text}.

\bibitem[Battisti \& Moras 2016]{battistimoras2016}
Battisti, Elisa, and Viviane Tebaldi Moras. 2016.
A vocalização da consoante lateral em coda silábica em uma variedade de
português brasileiro: análise sociolinguística em tempo real.
\textit{Gragoatá} 21(40): 90--112.
\href{https://doi.org/10.22409/gragoata.v21i40.33376}{doi:10.22409/gragoata.v21i40.33376}.

\bibitem[Dryer 2013]{dryer2013neg}
Dryer, Matthew S. 2013.
Order of negative morpheme and verb.
In Matthew S. Dryer and Martin Haspelmath (eds.),
\textit{The World Atlas of Language Structures Online}, chapter 143.
Leipzig: Max Planck Institute for Evolutionary Anthropology.
\url{https://wals.info/chapter/143}.

\bibitem[Embick \& Noyer 2007]{embicknoyer2007}
Embick, David, and Rolf Noyer. 2007.
Distributed Morphology and the syntax/morphology interface.
In Gillian Ramchand and Charles Reiss (eds.),
\textit{The Oxford Handbook of Linguistic Interfaces}, 289--324.
Oxford: Oxford University Press.
\href{https://doi.org/10.1093/oxfordhb/9780199247455.013.0010}{doi:10.1093/oxfordhb/9780199247455.013.0010}.

\bibitem[Gouskova 2012]{gouskova2012}
Gouskova, Maria. 2012.
Unexceptional segments.
\textit{Natural Language \& Linguistic Theory} 30(1): 79--133.
\href{https://www.gouskova.com/wp-content/uploads/2016/03/gouskova_2012_nllt.pdf}{Author's manuscript}.

\bibitem[Hall 2006]{hall2006}
Hall, Nancy. 2006.
Cross-linguistic patterns of vowel intrusion.
\textit{Phonology} 23(3): 387--429.
\href{https://doi.org/10.1017/S0952675706000996}{doi:10.1017/S0952675706000996}.

\bibitem[Han 2001]{han2001}
Han, Chung-hye. 2001.
Force, negation and imperatives.
\textit{The Linguistic Review} 18(4): 289--325.
\href{https://doi.org/10.1515/tlir.2001.004}{doi:10.1515/tlir.2001.004}.

\bibitem[Harley 2011]{harley2011}
Harley, Heidi. 2011.
Affixation and the Mirror Principle.
In Raffaella Folli and Christiane Ulbrich (eds.),
\textit{Interfaces in Linguistics: New Research Perspectives}, 166--186.
Oxford: Oxford University Press.
\href{https://heidiharley.com/pubs/affixation-and-the-mirror-principle/}{Author's publication page}.

\bibitem[Scheer 2017]{scheer2017}
Scheer, Tobias. 2017.
Voice-induced vowel lengthening.
\textit{Papers in Historical Phonology} 2: 116--151.
\href{https://doi.org/10.2218/pihph.2.2017.1910}{doi:10.2218/pihph.2.2017.1910}.

\bibitem[Unuk 2023]{unuk2023}
Unuk, Drago. 2023.
Fonetično-fonološki status fonema /v/ v slovenskem govorjenem knjižnem jeziku.
\textit{Jezikoslovni zapiski} 29(2): 37--75.
\href{https://doi.org/10.3986/JZ.29.2.02}{doi:10.3986/JZ.29.2.02}.

\bibitem[Veselinova 2016]{veselinova2016}
Veselinova, Ljuba N. 2016.
The negative existential cycle viewed through the lens of comparative data.
In Elly van Gelderen (ed.), \textit{Cyclical Change Continued}, 139--188.
Linguistik Aktuell/Linguistics Today 227.
Amsterdam: John Benjamins.
\href{https://doi.org/10.1075/la.227.06ves}{doi:10.1075/la.227.06ves}.

\bibitem[Walker 2001]{walker2001}
Walker, Rachel. 2001.
Positional markedness in vowel harmony.
In Caroline F\'ery, Antony Dubach Green, and Ruben van de Vijver (eds.),
\textit{Proceedings of HILP5}, 212--232.
Linguistics in Potsdam 12. Potsdam: University of Potsdam.
\href{https://bpb-us-e1.wpmucdn.com/sites.ucsc.edu/dist/f/1313/files/2021/10/HILP5.pdf}{Author's manuscript}.

\end{thebibliography}
\endgroup
